\documentclass[11pt]{article}
\usepackage[margin=0.9in]{geometry}
\usepackage{amsmath,amssymb,amsthm,booktabs,mathrsfs}
\usepackage[T1]{fontenc}
\usepackage{lmodern,microtype}
\usepackage[colorlinks=true,linkcolor=blue,urlcolor=blue]{hyperref}
\newcommand{\norm}[1]{\left\lVert#1\right\rVert}
\newtheorem{theorem}{Theorem}[section]
\newtheorem{lemma}[theorem]{Lemma}
\newtheorem{proposition}[theorem]{Proposition}
\newtheorem{corollary}[theorem]{Corollary}
\theoremstyle{remark}
\newtheorem{remark}[theorem]{Remark}
\title{Toward a four-dimensional Yang--Mills construction\\
\large Notebook 38: positive full-measure embeddings and local defect bounds}
\author{Research proof draft}
\date{17 October 2026}
\begin{document}
\maketitle
\begin{abstract}
We move beyond the homogeneous local integral to exact statements
about finite positive Wilson measures. A specified measurable
completion of the actual local coarse map admits positive linear
good and remainder operators whose sum is the exact pushforward.
Their iterates retain every history and the full normalization.
Under literal input-domain containment and the established local
Ward identities, the good measure has the actual local analytic
germ on each fixed finite regulator.

Independently, a site-plane slab argument bounds prescribed
large-curvature cells under the normalized Wilson measure, uniformly
in the total volume. A compact-group conditional Laplace theorem
controls fluctuations in a fixed box with small boundary data.
A boundary-only tree gauge joins these results: prescribed
separated boxes have a joint exponentially small probability of
failure of explicit local gauge/minimizer windows. These are
estimates on the entire bare Wilson measure, with the denominator
included.

The local box windows have not been identified with the averaged
constraint windows of the iterated renormalization map. The
all-history decomposition has not been bounded in a multiscale
complete-density norm. We record this remaining interface, a
precise match to the published large-field induction, and a
regularity counterexample ruling out one insufficient shortcut.
The continuum existence theorem is not established.
\end{abstract}

\section*{Relation to the construction and independence of inputs}

Notebook 37 establishes a positive logarithmic drift and arbitrarily
long finite weak homogeneous local histories within the working
proof chain and its explicitly stated upstream hypotheses. Those
histories alone do not represent the complete Wilson measure.
Here we address that next interface.

There are two distinct kinds of result below. The positive
embedding uses the previously defined ordinary composite coarse
map, stationary chart, pivot measure and local gauge Ward identity.
Its local-germ assertion is conditional on literal containment of
the real integration window in the represented input density.
By contrast, the bare Wilson probability results are proved
directly from compact-group integration, finite lattice geometry,
and positivity. They do not require the earlier multiscale
determinant calculation.

All field and coupling arguments in the new probability bounds
are real. The compact gauge groups and link Haar measures are
normalized to mass one. Unless explicitly stated otherwise, a
fixed-box constant may depend on its side length and the group,
but not on the surrounding torus volume. Such a constant is not
claimed uniform in a growing physical box at vanishing lattice
spacing. The general \(SU(N)\) bare-measure lemmas do not extend
all earlier \(SU(2)\) construction inputs to arbitrary groups.

The source repository remains pinned at
\texttt{e392ef45013a581664bab76b90867170b30fe10e}. This document is
a mathematical research draft with separate internal audits, not
an assertion of external verification of the preceding notebooks.
The exact status of each result and each missing implication is
part of its statement.

